\chapter{Computational Hardness and One-Way Functions}

In Chapter~7, we used probability and statistical distance to measure uncertainty. If we demand that a ciphertext leak zero information even to an adversary with unlimited time and computers, Shannon's theorem (which we prove next chapter) dictates an impossible price: the secret key must be as long as everything you will ever send.

Modern cryptography escapes this trap by changing the question. Instead of asking ``can the adversary \emph{ever} recover the message?'' we ask ``can the adversary recover the message \emph{using a realistic amount of computation}?'' In exchange for bounding the adversary's computation, we get short keys, public keys, and digital signatures. That trade rests on one bet: that certain problems are \emph{easy to pose and hard to solve}. Multiplying two large primes is fast; recovering them from the product seems to require effort that grows wildly with their size. This chapter builds the grammar of computational security -- polynomial time, $\PPT$ adversaries, negligible functions, one-way functions, trapdoors, and reductions -- used in every security proof in the rest of the book.

\chapterline

\section{The asymmetry we are buying}
\label{ch08:asymmetry}

Picture the usual cast. Alice wants to send Bob something secret; Eve listens to everything that crosses the wire. In the unconditional world Eve has unlimited computing power. In the computational world Eve is merely the best-funded adversary you can imagine -- a nation-state with warehouses of hardware -- but still bound by the same laws of computation as everyone else. Our schemes are allowed to be breakable in principle, as long as breaking them would take Eve longer than she (or the universe) has.

For that to be a sensible bargain we need a job that is cheap for Alice and Bob and ruinously expensive for Eve. The cleanest example is a padlock. Snapping a padlock shut takes a second and no key. Opening it without the key takes a locksmith, a drill, and time. The lock is \emph{asymmetric}: the forward direction (lock it) is easy, the backward direction (open it without the key) is hard. Cryptography relies on mathematical versions of this padlock.

\begin{intuitionbox}
Three small concrete asymmetries, all of which we will use later:
\begin{itemize}
\item \textbf{Multiplication vs.\ factoring.} Given $p=61$ and $q=53$, computing $N = pq = 3233$ is instant. Given only $N=3233$, finding the factors takes real searching. For 2048-bit $N$ it is, as far as anyone knows, hopeless.
\item \textbf{Exponentiation vs.\ discrete log.} Working modulo the prime $p=23$, computing $5^{6} \bmod 23 = 8$ is a handful of multiplications. Going backwards -- given that $5^{x}\equiv 8$, find $x$ -- means searching for the exponent. With a 256-bit modulus and the right group, the backward search is the discrete logarithm problem, and it is believed hard.
\item \textbf{Hashing vs.\ inverting a hash.} A good hash $\Hash$ turns ``hello'' into a fixed 256-bit digest in microseconds. Given the digest alone, producing \emph{any} input that hashes to it should require trying astronomically many candidates.
\end{itemize}
In each case one direction is cheap and the other is expensive.
\end{intuitionbox}

Everything depends on what we mean by ``hard.'' A child can fail to factor $3233$; that does not make factoring hard in any useful sense. We need a definition of hardness that (a) talks about all input sizes at once rather than one number, (b) measures effort against the size of the input, and (c) survives the fact that a lucky adversary might break one instance now and then. Building that definition carefully is the work of Sections~\ref{ch08:ppt}--\ref{ch08:negligible}. With it in hand we can finally say what a one-way function is and prove things about it.

\begin{historybox}
The idea that security could rest on computational difficulty rather than perfect secrecy is due to a cluster of work in the 1970s. Whitfield Diffie and Martin Hellman at Stanford, building on independent work by Ralph Merkle, framed it openly in their 1976 paper ``New Directions in Cryptography'': they proposed building secrecy on functions that are easy to compute and infeasible to invert, and coined the modern notion of a \emph{trapdoor} one-way function. Within a year Ron Rivest, Adi Shamir, and Leonard Adleman at MIT turned that abstract wish into RSA, an actual trapdoor built from factoring. Despite attempts by intelligence agencies to restrict academic publication, this conceptual move -- from ``unbreakable'' to ``infeasible to break in practice'' -- let cryptography leave the domain of military one-time pads and become an open science. Yet the foundational question remains open: proving that one-way functions exist would settle $\mathrm{P} \ne \mathrm{NP}$.
\end{historybox}

\section{Measuring effort: input size and asymptotics}
\label{ch08:size}

When we say factoring ``gets hard for large numbers,'' large compared to what? The yardstick is not the numerical value $N$ but the number of bits it takes to \emph{write down} $N$, which is about $n = \log_2 N$. A 2048-bit modulus is a number with roughly 617 decimal digits; trying to divide it by every integer up to $\sqrt N$ would mean about $2^{1024}$ trial divisions. The \emph{input} is 2048 bits -- tiny, fits in a tweet -- yet the naive attack is exponential in that tiny input.

This is exactly the gap we want. We measure an algorithm's running time as a function $T(n)$ of the input length $n$ in bits, and we care about how $T$ grows.

\begin{definition}[Polynomial and exponential time]
\label{ch08:def-polytime}
An algorithm runs in \emph{polynomial time} if there is a constant $c$ and a bound $n_0$ such that on every input of length $n \ge n_0$ it halts within $T(n) \le n^{c}$ steps. We write $T(n) = n^{O(1)}$. An algorithm is \emph{exponential time} if its running time grows like $2^{\Omega(n^{\delta})}$ for some $\delta > 0$ -- it eventually outpaces every polynomial.
\end{definition}

Polynomial time is our stand-in for ``feasible.'' It is a generous, slightly leaky abstraction: an algorithm that takes $n^{50}$ steps is polynomial and yet useless in practice, while one that takes $1.000001^{n}$ is exponential and yet fine for the sizes you meet. But the abstraction is robust in a way that any fixed cutoff (``feasible means under $2^{80}$ steps'') is not. Polynomials are closed under addition, multiplication, and composition, so if you build a feasible procedure out of feasibly many feasible pieces, the whole thing stays feasible. That closure is what makes proofs go through, and it is why the entire theory is phrased asymptotically.

\begin{intuitionbox}
Why bits and not the number itself? Because the difficulty of a problem should be compared against how much it costs even to \emph{state} the problem. Writing down a 2048-bit number costs 2048 bits; an attack that costs $2^{2048}$ is exponentially more expensive than the question. If we measured against the value $N$ instead, trial division (about $\sqrt N$ steps) would look ``polynomial in $N$'' and we would fool ourselves into thinking factoring is easy. The input-length yardstick keeps us out of that trap.
\end{intuitionbox}

\begin{connectionbox}
The companion algorithms book develops this scaling in detail: why multiplication of $n$-bit numbers is polynomial (schoolbook is $O(n^2)$, and fast methods do better), why fast modular exponentiation $a^b \bmod N$ costs only $O(n)$ multiplications via repeated squaring rather than $b$ multiplications, and why the best known factoring algorithm (the general number field sieve) runs in time roughly $\exp\!\big(c\,n^{1/3}(\log n)^{2/3}\big)$ -- \emph{sub}exponential, but still far above polynomial. Those cost analyses are the empirical evidence behind every hardness assumption in this chapter.
\end{connectionbox}

\begin{figure}[H]
\centering
\begin{tikzpicture}
  \begin{axis}[
    width=11.8cm,
    height=6.2cm,
    xmin=0, xmax=80,
    ymin=0, ymax=100,
    axis lines=left,
    xlabel={Input length $n$ (bits)},
    ylabel={Computational effort (log scale / arbitrary units)},
    xlabel style={font=\small},
    ylabel style={font=\small},
    tick label style={font=\footnotesize},
    xtick={0,20,40,60,80},
    ytick={0,25,50,75,100},
    yticklabels={$1$, $2^{25}$, $2^{50}$, $2^{75}$, $2^{100}$},
    grid=major,
    grid style={dashed, gray!20},
    legend style={at={(0.03,0.95)}, anchor=north west, font=\footnotesize, fill=white, fill opacity=0.9, draw=gray!40},
    title style={font=\small},
    title={Asymptotic scaling: polynomial vs.\ sub-exponential vs.\ exponential}
  ]
    \addplot[domain=2:80, samples=80, line width=1.5pt, sourceblue] {3 * (ln(x)/ln(2))};
    \addlegendentry{Polynomial $O(n^3)$ (Multiplication, AES, Eval)}
    \addplot[domain=2:80, samples=80, line width=1.5pt, quotientorange, dashed] {4.2 * (x^(1/3)) * ((ln(x)/ln(2))^(2/3))};
    \addlegendentry{Sub-exponential GNFS $\exp(c\,n^{1/3}\log^{2/3}\!n)$}
    \addplot[domain=2:80, samples=80, line width=1.8pt, trapred] {1.2 * x};
    \addlegendentry{Exponential $2^{\Omega(n)}$ (Brute-force key search)}
    \draw[dashed, structurepurple, line width=1pt] (axis cs:0,40) -- (axis cs:80,40)
      node[below, pos=0.72, font=\footnotesize\bfseries, text=structurepurple, fill=white, inner sep=1.5pt] {Practical limit ($\approx 2^{40}$ ops/sec)};
  \end{axis}
\end{tikzpicture}
\caption{The computational hardness gap. While polynomial-time forward operations (blue) remain tiny and flat even as keys grow, exponential brute-force attacks (red) skyrocket beyond reach. Sub-exponential attacks (orange dashed, such as the General Number Field Sieve for factoring) grow faster than any polynomial, explaining why RSA requires 3072-bit keys while 256-bit elliptic curves suffice.}
\label{ch08:fig-complexity-growth}
\end{figure}
\section{The adversary: probabilistic polynomial time}
\label{ch08:ppt}
We do not model Eve as one fixed program but as the \emph{best} program of her size. We model her as an algorithm, and we grant her two things beyond a plain deterministic computer.
First, randomness. Real attackers flip coins -- they pick random points to evaluate, guess random candidates, run randomized number-theory algorithms. Forbidding randomness would understate Eve's power and let us claim false security. So the adversary may toss fair coins for free.
Second, we hold her to polynomial time, because an adversary allowed exponential time can break essentially everything (she can brute-force any key). The interesting and achievable goal is security against \emph{efficient} adversaries.
Putting these together gives the standard model of a realistic attacker.
\begin{definition}[$\PPT$ adversary]
\label{ch08:def-ppt}
A \emph{probabilistic polynomial-time} ($\PPT$) algorithm is a randomized algorithm whose running time is bounded by a polynomial in the length of its input, on every input and every setting of its internal coins. A $\PPT$ \emph{adversary} is such an algorithm whose job is to break a scheme: recover a key, distinguish two ciphertexts, forge a signature, invert a function.
\end{definition}
Because Eve gets coins, the outcome of an attack is a random variable. ``Eve breaks the scheme'' is an \emph{event with a probability}, not a yes/no fact. This is why the next ingredient -- a way to talk about probabilities that are ``so small they do not count'' -- is central to the framework.
\begin{focusbox}
We will attach a \emph{security parameter}, usually called $n$ (or $\lambda$ in other texts, or $k$), to every scheme. Think of $n$ as a dial: turn it up and keys get longer, the legitimate work grows polynomially, and -- if the design is sound -- the adversary's success probability shrinks faster than any polynomial. Key length is roughly $n$; an RSA modulus of $n$ bits is the security parameter made concrete. Every quantity below is a \emph{function of $n$}: the adversary's run time, her advantage, the time to encrypt. Holding $n$ in mind as ``the knob'' is the single most useful habit for reading this material.
\end{focusbox}
\section{Negligible vs.\ noticeable}
\label{ch08:negligible}
Suppose an attack against a 256-bit key succeeds with probability $2^{-128}$. That is no cause for worry: even running the attack a trillion times a second since the Big Bang, the expected number of successes is far below one. A probability that small is, for every practical and theoretical purpose, zero. We need to make ``that small'' precise -- and make it scale with the security parameter, so that turning up $n$ actually buys safety.
\begin{definition}[Negligible function]
\label{ch08:def-negl}
A function $\negl : \N \to \R_{\ge 0}$ is \emph{negligible} if for every positive constant $c$ there is a bound $n_0$ such that
\[
\negl(n) < \frac{1}{n^{c}} \qquad \text{for all } n \ge n_0 .
\]
Equivalently: $\negl$ eventually drops below the reciprocal of \emph{every} polynomial.
\end{definition}
The definition says ``smaller than $1/n^{c}$ for all $c$,'' which sounds abstract until you plug in examples. The functions $2^{-n}$, $2^{-\sqrt n}$, and $n^{-\log n}$ are all negligible: each one eventually undercuts $1/n^{1000}$, $1/n^{10^6}$, and so on. By contrast $1/n^{3}$, $1/(1000\,n^2)$, and $1/2^{\log_2 n} = 1/n$ are \emph{not} negligible -- they are polynomially large. A function that is not negligible (more precisely, that exceeds $1/n^c$ for some fixed $c$ infinitely often) is called \emph{noticeable}.
\begin{example}[Negligible or not?]
\label{ch08:ex-negl}
Sort these by plugging in the definition.
\begin{itemize}
\item $f(n) = 2^{-n}$: at $n=80$ this is about $10^{-24}$; it beats $1/n^c$ for every $c$ once $n$ is large enough. \textbf{Negligible.}
\item $f(n) = 2^{-n/2}$: halving the exponent only changes the constant in front; still beats every polynomial. \textbf{Negligible.} (This is why ``128-bit security'' against a $2^{n/2}$ birthday attack wants $n = 256$.)
\item $f(n) = 1/n^{100}$: matches the polynomial $1/n^{100}$, so it cannot stay below $1/n^{101}$. \textbf{Noticeable.}
\item $f(n) = n^{-\log n}$: take logs, $\log f = -(\log n)^2$, which beats $-c\log n$ for every $c$. \textbf{Negligible.}
\item $f(n) = 2^{-n} + 1/n^{5}$: a sum is negligible only if both parts are; the $1/n^5$ term dominates. \textbf{Noticeable.}
\end{itemize}
\end{example}
It helps to watch the crossover happen with actual numbers. The negligible function $2^{-n}$ starts out \emph{larger} than the polynomial $1/n^{3}$ for small $n$ and only later dives underneath it -- ``eventually'' is doing real work in the definition.
\begin{center}
\begin{tabular}{r r r l}
\toprule
$n$ & $2^{-n}$ & $1/n^{3}$ & which is smaller? \\
\midrule
$5$    & $3.1\times 10^{-2}$  & $8.0\times 10^{-3}$  & $1/n^3$ (poly wins) \\
$10$   & $9.8\times 10^{-4}$  & $1.0\times 10^{-3}$  & roughly equal (crossover) \\
$20$   & $9.5\times 10^{-7}$  & $1.3\times 10^{-4}$  & $2^{-n}$ \\
$40$   & $9.1\times 10^{-13}$ & $1.6\times 10^{-5}$  & $2^{-n}$ \\
$80$   & $8.3\times 10^{-25}$ & $2.0\times 10^{-6}$  & $2^{-n}$ \\
\bottomrule
\end{tabular}
\end{center}
Past the crossover near $n=10$ the exponential plunges and never resurfaces; by $n=80$ it is twenty orders of magnitude below the polynomial, and the gap only widens. Negligibility is the claim that this crossing happens against \emph{every} polynomial, no matter how large its degree -- a higher-degree polynomial just pushes the crossover point to a larger $n$.
We use this exact definition because it behaves well under the two kinds of composition that proofs rely on.
\begin{whybox}
\textbf{Negligible $\times$ polynomial is negligible.} If $\negl(n)$ is negligible and $p(n)$ is any polynomial, then $p(n)\cdot\negl(n)$ is still negligible. \emph{Proof sketch.} Fix any constant $c$. We must show $p(n)\,\negl(n) < 1/n^{c}$ eventually. Since $p$ is a polynomial, $p(n) \le n^{d}$ for some $d$ and all large $n$. Because $\negl$ is negligible, it eventually falls below $1/n^{c+d}$. Multiply: $p(n)\,\negl(n) \le n^{d}\cdot n^{-(c+d)} = n^{-c}$, as needed. $\square$
This is the property we lean on constantly. A $\PPT$ adversary runs polynomially many steps; if each step or each query succeeds only with negligible probability, a union bound over those steps multiplies a polynomial by a negligible, and the total stays negligible. An attacker cannot amplify a negligible edge into a real one by repeating a polynomial number of times.
\textbf{Negligible $+$ negligible is negligible.} The sum of two negligible functions is negligible (each is below $\tfrac12 n^{-c}$ eventually, so the sum is below $n^{-c}$). So a proof that strings together polynomially many negligible failure events still concludes ``negligible.''
\end{whybox}
These closure facts are exactly why we did \emph{not} pick a fixed threshold like ``probability under $2^{-80}$.'' A fixed threshold is not closed under the operations proofs perform -- run a $2^{-80}$ attack $2^{40}$ times and you cross it. The asymptotic definition is engineered so that the natural manipulations of a security proof (repeat a polynomial number of times, union-bound over polynomially many cases, compose two reductions) all preserve negligibility.
\begin{trapbox}
Negligible is an \emph{asymptotic} statement about a function of $n$, not a property of a single number. ``Is $2^{-100}$ negligible?'' is a category error -- $2^{-100}$ is just a number. The right question is whether the \emph{family} of success probabilities, as you scale $n$, eventually beats every polynomial. In engineering we do fix a concrete $n$ (``use a 3072-bit modulus'') and then talk about a concrete number like $2^{-128}$; that is the asymptotic theory evaluated at one operating point. Keep the two registers straight: theory quantifies over $n$, deployment picks an $n$.
\end{trapbox}
\section{Advantage: how much better than guessing}
\label{ch08:advantage}
Many attacks do not recover a key outright; they merely \emph{distinguish}. Eve sees a ciphertext and tries to tell whether it encrypts ``attack at dawn'' or ``attack at noon,'' or to tell a real cipher output from random noise. Even an adversary who ignores the ciphertext and flips a coin wins half the time. So raw success probability overstates her skill. The meaningful quantity is how far above the trivial baseline she climbs.
\begin{definition}[Distinguishing advantage]
\label{ch08:def-adv}
Let $\mathcal D$ be a distinguisher (a $\PPT$ algorithm outputting a single bit) faced with one of two experiments, $\mathsf{Exp}_0$ and $\mathsf{Exp}_1$, chosen with probability $\tfrac12$ each. Its \emph{advantage} is
\[
\Adv(\mathcal D) \;=\; \Bigl| \Pr[\mathcal D = 1 \mid \mathsf{Exp}_1] \;-\; \Pr[\mathcal D = 1 \mid \mathsf{Exp}_0] \Bigr| .
\]
Equivalently, if the experiment is chosen by a hidden fair coin $b$ and $\mathcal D$ outputs a guess $b'$, then $\Pr[b' = b] = \tfrac12 + \tfrac12\Adv(\mathcal D)$.
\end{definition}
The two forms say the same thing. Guessing blindly gives $\Pr[b'=b]=\tfrac12$ and advantage $0$. A perfect distinguisher gives $\Pr[b'=b]=1$ and advantage $1$. Advantage is the slice of certainty Eve adds on top of coin-flipping, rescaled to live in $[0,1]$.
\begin{example}[A tiny advantage calculation]
\label{ch08:ex-adv}
Suppose a flawed cipher leaks the last bit of the plaintext into the ciphertext one time in eight; the rest of the time the ciphertext looks uniform. Eve's strategy: if she spots the leak, output the leaked bit as her guess of $b$; otherwise flip a coin. With probability $\tfrac18$ she sees the leak and is always right; with probability $\tfrac78$ she guesses and is right half the time:
\[
\Pr[b'=b] \;=\; \tfrac18\cdot 1 \;+\; \tfrac78\cdot\tfrac12 \;=\; \tfrac{1}{8} + \tfrac{7}{16} \;=\; \tfrac{9}{16}.
\]
So $\tfrac12 + \tfrac12\Adv = \tfrac{9}{16}$, giving $\Adv = \tfrac18$. A non-negligible advantage of one eighth means the cipher is broken: Eve does measurably better than chance, and that is all our security definitions forbid.
\end{example}
\begin{focusbox}
\textbf{What ``secure'' means, in one sentence.} A scheme is secure (in the computational sense) when \emph{every} $\PPT$ adversary has only \emph{negligible} advantage, as a function of the security parameter $n$. Spelled out: for all $\PPT$ $\mathcal A$ there is a negligible $\negl$ with $\Adv(\mathcal A) \le \negl(n)$. Read it as a contest with three quantifiers -- \emph{for all} efficient strategies, the edge over guessing is \emph{negligible} in $n$. Every formal definition in later chapters (semantic security, CPA security, unforgeability, collision resistance) is this template with a specific experiment plugged in.
\end{focusbox}
\section{One-way functions}
\label{ch08:owf}
Now we can name the central object. A one-way function is the mathematical padlock: easy to apply, hard to undo. The two ``easy'' and ``hard'' clauses are exactly the polynomial-time and negligible-advantage notions we just built.
\begin{definition}[One-way function]
\label{ch08:def-owf}
A function $f : \{0,1\}^{*} \to \{0,1\}^{*}$ is \emph{one-way} if:
\begin{enumerate}
\item \textbf{Easy to compute.} There is a deterministic polynomial-time algorithm that, given $x$, outputs $f(x)$.
\item \textbf{Hard to invert.} For every $\PPT$ algorithm $\mathcal A$, the probability that $\mathcal A$ produces \emph{some} preimage is negligible:
\[
\Pr_{x \getsr \{0,1\}^{n}}\bigl[\, \mathcal A\bigl(1^{n}, f(x)\bigr) \in f^{-1}\bigl(f(x)\bigr)\,\bigr] \;\le\; \negl(n).
\]
\end{enumerate}
The probability is over the random choice of $x$ \emph{and} over $\mathcal A$'s coins.
\end{definition}
Several details in that definition deserve attention, because each one closes a loophole that would otherwise make the notion vacuous or immediately false.
\begin{walkbox}
\emph{Why ``some preimage'' and not ``the same $x$''?} Because $f$ need not be injective. If $f(x)=f(x')$, any inverter that outputs $x'$ has done its job -- it found something mapping to the target. Demanding the original $x$ would be unfair when the information to distinguish $x$ from $x'$ is simply not present in $f(x)$.
\emph{Why is $x$ chosen at random?} One-wayness is an \emph{average-case} property. We are not claiming $f$ is hard to invert on its worst input; we are claiming it is hard on a \emph{random} input. That is a stronger, more useful demand for cryptography, where keys and messages are random. A function could be hard in the worst case yet trivial on the inputs you actually meet -- useless to us.
\emph{What is that $1^{n}$ doing as an argument?} It is the security parameter written in \emph{unary} -- a string of $n$ ones. Handing $\mathcal A$ a string of length $n$ guarantees that ``polynomial time'' for $\mathcal A$ means polynomial in $n$, even when $f(x)$ happens to be short. Without it, an adversary could be ``efficient'' relative to a tiny input yet effectively be granted very little time. It is a bookkeeping device that pins down the clock.
\end{walkbox}
\begin{whybox}
\textbf{Why both clauses are mandatory.} Drop ``easy to compute'' and you can declare any random-looking map one-way without being able to use it -- legitimate parties could not even apply the function. Drop ``hard to invert'' and you have an ordinary function, no padlock at all. The notion is useful only because it combines both: useful (fast forward) \emph{and} safe (slow backward). A constant function $f(x)=0$ is immediately easy to compute but immediately easy to invert (output anything); the identity $f(x)=x$ is easy both ways. One-wayness lives in the narrow band where forward is cheap and backward is hopeless on average.
\end{whybox}
\begin{intuitionbox}
A homely model: dropping a wine glass. Going from glass to shards is one function evaluation and takes a flick of the wrist. Going from shards back to glass is, thermodynamically, possible -- no law forbids it -- but the number of shard arrangements you would have to search through makes it hopeless in practice. One-way functions are the digital version: the forward map mixes information so thoroughly that running it backward means searching a space too large to search.
\end{intuitionbox}
\subsection{Candidate one-way functions}
\label{ch08:candidates}
\emph{Nobody has proved that one-way functions exist.} Their existence would imply $\mathrm P \ne \mathrm{NP}$, one of the great open problems. What we have instead are \emph{candidates} -- functions that have resisted decades of expert attack and that we \emph{believe} are one-way. Cryptography is, in this sense, an empirical science resting on unproven (but heavily stress-tested) hardness conjectures.
\begin{example}[Multiplication / the factoring candidate]
\label{ch08:ex-mult}
Define $f(p,q) = pq$ where $p,q$ are $n$-bit primes (this is the multiplication function restricted to primes). Forward: multiply, polynomial time. Backward: given $N=pq$, recover $\{p,q\}$ -- this is integer factoring. With $p=61, q=53$ we get $N=3233$; inverting means factoring $3233$. At 2048 bits, inverting is the problem RSA's security leans on.
\end{example}
\begin{example}[Modular exponentiation / the discrete-log candidate]
\label{ch08:ex-exp}
Fix a large prime $p$ and a generator $g$ of $\Zns{p}$. Define $f(x) = g^{x} \bmod p$. Forward: fast exponentiation, $O(n)$ multiplications. Backward: given $y=g^x$, find $x$ -- the discrete logarithm. With $p=23$, $g=5$, and $x=6$ we get $y = 5^6 \bmod 23 = 8$; inverting means solving $5^{x}\equiv 8 \pmod{23}$ for $x$.
\end{example}
\begin{example}[A hash as a one-way function]
\label{ch08:ex-hash}
A cryptographic hash $\Hash:\{0,1\}^{*}\to\{0,1\}^{256}$ is designed to be \emph{preimage resistant}: given a digest $d$, finding any $m$ with $\Hash(m)=d$ should cost about $2^{256}$ work. Preimage resistance is precisely one-wayness for $\Hash$. (Hashes ask for more -- collision resistance too -- which we treat in Chapter~14.)
\end{example}
\begin{connectionbox}
The discrete-log toy in Example~\ref{ch08:ex-exp} is worth tracing by hand, because it shows the asymmetry in miniature. Forward, in $\Zns{23}$ with $g=5$: the powers cycle as $5^1{=}5,\;5^2{=}2,\;5^3{=}10,\;5^4{=}4,\;5^5{=}20,\;5^6{=}8,\dots$ (each step is ``multiply by 5, reduce mod 23''). Backward, you are handed $8$ and must find which exponent produced it -- here you would scan the list. For $p$ of a few hundred bits the list has more entries than there are atoms in the observable universe, and the best known general attacks (baby-step giant-step, Pollard's rho, index calculus) still cost far more than the polynomial forward direction. The number-theory appendix proves $5$ is a generator mod $23$ and explains why a generator makes the forward map a bijection onto $\Zns{p}$.
\end{connectionbox}
\begin{implementbox}
A toy ``one-way'' function and a brute-force inverter, to feel the asymmetry. This is for intuition only: real one-way functions use cryptographic sizes, and a 16-bit modulus like this one is broken in microseconds.
\begin{lstlisting}[language=Python]
# LEARNING ONLY -- toy sizes, NOT secure. Do not use for anything real.
p = 65537          # a small prime (2^16 + 1)
g = 3              # a generator of the multiplicative group mod p
def f(x):
    # forward direction: fast, O(log x) multiplications via pow()
    return pow(g, x, p)
def invert_bruteforce(y):
    # backward direction: search every exponent until one matches.
    # Cost grows with p -- here ~65k tries; at cryptographic sizes
    # this loop would outlast the universe.
    for x in range(1, p):
        if pow(g, x, p) == y:
            return x
    return None
secret = 41234
y = f(secret)                       # easy: one call
print("public value y =", y)
print("recovered x   =", invert_bruteforce(y))   # "hard": brute force
# Notice: f(secret) returned instantly. invert_bruteforce had to grind
# through up to p candidates. That gap -- tiny here -- is the entire
# bet of modern cryptography, blown up to 2048+ bits.
\end{lstlisting}
What matters here is the shape of the cost: \texttt{f} is one operation, \texttt{invert\_bruteforce} is a loop over the whole space. Scale $p$ to thousands of bits and the loop becomes physically impossible while \texttt{f} stays instant.
\end{implementbox}
The brute-force loop above is the simplest possible inverter; it scans all $p$ candidates. The real reason discrete log is only \emph{conjectured} hard -- not provably $2^{n}$ hard -- is that cleverer attacks shave the exponent. The simplest of these is worth running by hand, because it shows both that the forward/backward gap is real \emph{and} that it is smaller than naive counting suggests.
\begin{labbox}
\textbf{Baby-step giant-step on a tiny group.} We solve $5^{x}\equiv 9 \pmod{23}$ -- recover the discrete log of $9$ to base $5$ in $\Zns{23}$ -- in far fewer than $22$ steps. The group has order $m = 22$; set $t = \lceil\sqrt{22}\rceil = 5$. Write the unknown as $x = it + j$ with $0 \le j < t$ and $0 \le i \le t$.
\emph{Baby steps.} Tabulate $5^{j}\bmod 23$ for $j = 0,\dots,4$:
\[
5^{0}{=}1,\quad 5^{1}{=}5,\quad 5^{2}{=}2,\quad 5^{3}{=}10,\quad 5^{4}{=}4.
\]
Store the pairs $\{(1,0),(5,1),(2,2),(10,3),(4,4)\}$ keyed by the power.
\emph{Giant steps.} The giant stride is $5^{-t} = 5^{-5}\bmod 23$. Since $5^{5}\equiv 20$ and $20^{-1}\equiv 15 \pmod{23}$ (because $20\cdot 15 = 300 = 13\cdot 23 + 1$), the stride factor is $15$. Now compute $9\cdot 15^{i}\bmod 23$ for $i=0,1,2,\dots$ and look each result up in the baby table:
\[
i{=}0:\;9,\quad i{=}1:\;9\cdot 15 = 135 \equiv 20,\quad i{=}2:\;20\cdot 15 = 300 \equiv 1.
\]
At $i=2$ the value $1$ is in the table, paired with $j=0$. So $x = it + j = 2\cdot 5 + 0 = 10$.
\emph{Check.} $5^{10}\bmod 23$: $5^{5}\equiv 20$, so $5^{10}\equiv 20^{2} = 400 \equiv 9 \pmod{23}$. Correct.
We found $x=10$ in $5$ baby steps plus $3$ giant steps -- $8$ multiplications instead of scanning $22$ exponents. The win is $O(\sqrt m)$ versus $O(m)$, the same square-root that the connection box above attributes to the birthday bound. For a $256$-bit group $\sqrt m \approx 2^{128}$, still hopeless; but the $2\times$ in the exponent is exactly why ``128-bit security'' demands a $256$-bit group, not a $128$-bit one.
\end{labbox}
\section{Trapdoors: a hidden shortcut}
\label{ch08:trapdoor}
A one-way function locks the door for everyone, including the person who set it up. That is great for, say, password storage (store $\Hash(\text{password})$; nobody can run it backward). But it is hopeless for sending Bob a secret: if Bob cannot invert the function either, the message is lost to him too. We need a function that is hard to invert \emph{in general} but easy to invert for someone holding a secret key. That secret is the \emph{trapdoor}.
\begin{definition}[Trapdoor one-way function]
\label{ch08:def-tdf}
A \emph{trapdoor one-way function} is a family of functions $\{f_{pk}\}$ together with a generation algorithm $\Gen$ that, on input $1^{n}$, outputs a pair $(pk, td)$ -- a public index and a trapdoor -- such that:
\begin{enumerate}
\item Given $pk$, the function $f_{pk}$ is easy to compute.
\item Without the trapdoor, $f_{pk}$ is one-way: for every $\PPT$ $\mathcal A$, $\Pr[\mathcal A(pk, f_{pk}(x)) \in f_{pk}^{-1}(f_{pk}(x))] \le \negl(n)$.
\item \textbf{With} the trapdoor $td$, inversion is easy: there is a polynomial-time algorithm that, given $td$ and $y=f_{pk}(x)$, outputs a preimage of $y$.
\end{enumerate}
\end{definition}
The padlock analogy extends to this case. A plain one-way function is a padlock with no key at all -- you can snap it shut but nobody opens it. A trapdoor function is a padlock \emph{with} a key: anyone can click it shut (the public direction), but only the keyholder reopens it (the trapdoor). Handing out open padlocks while keeping the key is the idea behind public-key cryptography.
\begin{intuitionbox}
RSA, which gets all of Chapter~19, is the canonical trapdoor. The public index is $(N,e)$ with $N=pq$; the function is $f(x)=x^{e}\bmod N$, easy for anyone. The trapdoor is the factorization $p,q$, from which one computes the secret exponent $d$ that undoes the $e$-th power. Tiny instance: $p=61, q=53, N=3233$. Knowing $p,q$ lets you compute $d$ and decrypt in a flash; knowing only $N=3233$ leaves you factoring. Same asymmetry as before, now with a back door for the rightful owner.
\end{intuitionbox}
\begin{figure}[H]
\centering
\begin{tikzpicture}[
  font=\small,
  card/.style={draw=#1!70!black, fill=#1!8, rounded corners=3pt, line width=0.8pt, inner sep=4pt},
  lanebox/.style={draw=gray!30, fill=gray!2, rounded corners=4pt, dashed, line width=0.6pt},
  arrowstyle/.style={-{Stealth[scale=1.0]}, line width=1.1pt}
]
  % Public Forward Lane
  \node[card=sourceblue, minimum width=2.4cm, minimum height=1.1cm, align=center] (x) at (0, 2.3)
    {\textbf{Plaintext / Seed}\\[1pt]$x \in \{0,1\}^n$};
  \node[card=sourceblue, minimum width=3.2cm, minimum height=1.1cm, align=center] (f) at (5.6, 2.3)
    {\textbf{Public Function}\\[1pt]$y = f_{pk}(x)$};
  \node[card=targetgreen, minimum width=2.4cm, minimum height=1.1cm, align=center] (y) at (11.2, 2.3)
    {\textbf{Ciphertext / Image}\\[1pt]$y \in \{0,1\}^n$};

  \draw[arrowstyle, sourceblue] (x.east) -- node[above=2pt, font=\scriptsize\bfseries, text=sourceblue] {Easy (poly time)} (f.west);
  \draw[arrowstyle, sourceblue] (f.east) -- node[above=2pt, font=\scriptsize, text=gray!80!black] {Public key $pk$} (y.west);

  % Adversarial Barrier (Middle)
  \node[card=trapred, minimum width=6.2cm, minimum height=1.2cm, align=center] (wall) at (5.6, 0.4)
    {\textbf{Hard Inversion Barrier (No Trapdoor)}\\[1pt]
     \footnotesize Adversary sees $(pk, y)$ only $\implies$ Brute force search $2^{\Omega(n)}$\\[1pt]
     \scriptsize Inversion success probability $\le \negl(n)$ for all $\PPT$ adversaries};

  \draw[arrowstyle, trapred, dashed] ([yshift=-2mm]y.south) |- (wall.east);
  \draw[arrowstyle, trapred, dashed] (wall.west) -| ([yshift=-2mm]x.south);
  \node[font=\footnotesize\bfseries, text=trapred] at (5.6, -0.6) {\texttimes{} Infeasible for $\PPT$ algorithms};

  % Private Trapdoor Lane (Bottom)
  \node[card=quotientorange, minimum width=5.6cm, minimum height=1.2cm, align=center] (tdbox) at (5.6, -2.0)
    {\textbf{Trapdoor Shortcut ($td$ holder)}\\[1pt]
     \footnotesize Trapdoor inversion algorithm: $x = f^{-1}(y \mid td)$\\[1pt]
     \scriptsize Instant polynomial-time recovery using private key};

  \draw[arrowstyle, quotientorange] ([xshift=4mm]y.south) |- (tdbox.east);
  \draw[arrowstyle, quotientorange] (tdbox.west) -| ([xshift=-4mm]x.south);
  \node[font=\footnotesize\bfseries, text=quotientorange] at (5.6, -3.0) {Easy with private trapdoor $td$ (e.g.\ factoring $p,q$ or discrete log $x$)};

  % Background container
  \begin{scope}[on background layer]
    \node[lanebox, fit={(x) (y) (f) (wall) (tdbox) (5.6,-3.2)}, inner sep=8pt] {};
  \end{scope}
\end{tikzpicture}
\caption{The three-way architecture of a trapdoor one-way function. Moving forward (blue) is computationally easy for anyone holding public key $pk$. Moving backward without the secret (red dashed) requires exhaustive searching over an exponentially large space. But holding the secret trapdoor $td$ (orange) provides an algebraic shortcut that inverts the function in polynomial time.}
\label{ch08:fig-trapdoor}
\end{figure}

\begin{trapbox}
``Hard to invert without the trapdoor'' must hold even though $pk$ is \emph{public}. The adversary sees the full description of $f_{pk}$ -- the modulus, the exponent, everything except $td$. Security never rests on hiding the algorithm or the public key; it rests entirely on the secrecy of $td$ and the conjectured hardness of inverting. This is Kerckhoffs's principle in computational form: assume Eve knows the system, and make the secret the only thing she lacks.
\end{trapbox}

\section{Reductions: turning hardness into security}
\label{ch08:reductions}

We have hard problems (factoring, discrete log) and we want secure schemes. The bridge between them is a \emph{reduction}, and it is the single most important proof technique in the book. The logic runs by contraposition.

We want to prove: \emph{if problem $Y$ is hard, then scheme $X$ is secure.} Directly arguing ``no attack exists'' is hopeless -- you cannot enumerate all possible attacks. So we flip it: assume an attack exists and derive a contradiction.

\begin{focusbox}
\textbf{The shape of every reduction.} To prove ``scheme $X$ is secure assuming $Y$ is hard'':
\begin{enumerate}
\item Suppose, for contradiction, that some $\PPT$ adversary $\mathcal A$ breaks $X$ with non-negligible advantage $\eps(n)$.
\item Build a $\PPT$ algorithm $\mathcal B$ -- the \emph{reduction} -- that uses $\mathcal A$ as a subroutine (a ``black box'') to solve $Y$.
\item Show $\mathcal B$ solves $Y$ with probability related to $\eps(n)$, also non-negligible, in polynomial time.
\item Conclude: a $\PPT$ solver for $Y$ exists, contradicting ``$Y$ is hard.'' Hence no such $\mathcal A$ exists, i.e.\ $X$ is secure.
\end{enumerate}
The slogan: \emph{breaking $X$ lets you solve $Y$}; since $Y$ cannot be solved, $X$ cannot be broken.
\end{focusbox}

The hard part is step 2: $\mathcal B$ must feed $\mathcal A$ inputs that look exactly like a real attack scenario, then read off $\mathcal A$'s answer and \emph{translate} it into a solution to $Y$. $\mathcal B$ does not need to understand \emph{how} $\mathcal A$ works -- it treats $\mathcal A$ as an oracle that, when handed a valid-looking challenge, coughs up a break with probability $\eps(n)$.

\begin{figure}[H]
\centering
\begin{tikzpicture}[
  font=\small,
  card/.style={draw=#1!70!black, fill=#1!10, rounded corners=3pt, line width=0.8pt, inner sep=4pt},
  arrowstyle/.style={-{Stealth[scale=1.0]}, line width=1.1pt}
]
  % External Hard Problem Source
  \node[card=quotientorange, minimum width=2.8cm, minimum height=1.2cm, align=center] (inst) at (0, 0.9)
    {\textbf{Hard Instance $y \in Y$}\\[2pt]\footnotesize e.g.\ $N = pq$ (factoring)\\[1pt]\scriptsize or $(g, g^x)$ (discrete log)};

  % External Solution
  \node[card=targetgreen, minimum width=2.8cm, minimum height=1.2cm, align=center] (sol) at (0, -1.4)
    {\textbf{Solution to $Y$}\\[2pt]\footnotesize e.g.\ factors $p, q$\\[1pt]\scriptsize or exponent $x = \dlog_g(y)$};

  % Reduction Machine B
  \node[card=structurepurple, minimum width=7.2cm, minimum height=4.8cm] (B) at (7.0, -0.05) {};
  \node[font=\bfseries\small, text=structurepurple!80!black, anchor=north west] at ([xshift=6pt, yshift=-5pt]B.north west)
    {Reduction Algorithm $\mathcal{B}$ (Polynomial Time)};

  % Inside B: Simulation adapter
  \node[card=sourceblue, minimum width=2.8cm, minimum height=1.1cm, align=center, inner sep=4pt] (adapter) at (5.3, 0.9)
    {\textbf{Challenge Simulator}\\[3pt]\footnotesize Embeds $y$ into scheme $X$};

  % Inside B: Hypothetical Adversary A
  \node[card=trapred, minimum width=2.6cm, minimum height=1.2cm, align=center, inner sep=4pt] (A) at (9.0, -0.25)
    {\textbf{Adversary $\mathcal{A}$}\\[3pt]\footnotesize Breaks Scheme $X$};

  % Inside B: Extractor
  \node[card=targetgreen, minimum width=2.8cm, minimum height=1.1cm, align=center, inner sep=4pt] (extract) at (5.3, -1.4)
    {\textbf{Algebraic Extractor}\\[3pt]\footnotesize Solves hard problem $Y$};

  % Corridors & Flows
  % 1. Hard instance into B
  \draw[arrowstyle, quotientorange] (inst.east) -- node[above=2pt, font=\scriptsize, text=quotientorange!80!black] {Input $y$} (adapter.west);

  % 2. Adapter to A
  \draw[arrowstyle, sourceblue] (adapter.east) -| node[above=2pt, pos=0.32, font=\scriptsize, text=sourceblue!80!black] {Simulated $X$-challenge} (A.north);

  % 3. A returns break to Extractor
  \draw[arrowstyle, trapred] (A.south) |- node[below=2pt, pos=0.65, font=\scriptsize, text=trapred!80!black] {Attack break / leak} (extract.east);

  % 4. Extractor to Solution
  \draw[arrowstyle, targetgreen] (extract.west) -- node[above=2pt, font=\scriptsize\bfseries, text=targetgreen!80!black] {Solution} (sol.east);

  % Contradiction summary banner below
  \node[draw=gray!30, fill=gray!3, rounded corners=3pt, line width=0.6pt, align=center, font=\scriptsize, inner sep=5pt] at (4.4, -3.2)
    {\textbf{Proof by Contradiction:} If $\mathcal{A}$ runs in $\PPT$ with non-negligible advantage $\eps(n)$, then $\mathcal{B}$ solves hard problem $Y$ in $\PPT$\\
     with probability $\text{poly}(\eps(n))$. Since no efficient solver exists for $Y$, no such adversary $\mathcal{A}$ can exist: scheme $X$ is secure.};

\end{tikzpicture}
\caption{The mechanics of a cryptographic reduction. To prove scheme $X$ secure from hard problem $Y$, we build an efficient reduction $\mathcal{B}$ that wraps a hypothetical adversary $\mathcal{A}$. $\mathcal{B}$ feeds $\mathcal{A}$ a simulated environment derived from $Y$. When $\mathcal{A}$ breaks the scheme, $\mathcal{B}$ extracts a solution to $Y$, contradicting $Y$'s hardness.}
\label{ch08:fig-reduction}
\end{figure}

\begin{example}[A worked reduction: square roots are as hard as factoring]
\label{ch08:ex-rabin}
Let $N=pq$ with $p,q$ distinct odd primes, and consider the function $f(x)=x^{2}\bmod N$ (the basis of the Rabin trapdoor). We show: \emph{if you can compute square roots mod $N$, you can factor $N$.} So inverting $f$ is at least as hard as factoring -- a reduction \emph{from} factoring \emph{to} root-finding.

Suppose $\mathcal A(N,y)$ returns a square root of $y$ whenever one exists. Build $\mathcal B$ to factor $N$:
\begin{enumerate}
\item Pick a random $r$ in $\{1,\dots,N-1\}$ with $\gcd(r,N)=1$ (if the $\gcd$ is not $1$ you already found a factor -- a lucky shortcut).
\item Compute $y = r^{2}\bmod N$ and hand $(N,y)$ to $\mathcal A$. It returns some root $s$ with $s^{2}\equiv y \pmod N$.
\item Now $s^{2}\equiv r^{2}$, so $N \mid (s-r)(s+r)$. Because $N$ has two prime factors, $y$ has \emph{four} square roots mod $N$; $\mathcal A$ cannot tell which one $\mathcal B$ started from, so with probability $\tfrac12$ the returned $s$ is \emph{not} $\pm r$. In that case $\gcd(s-r,N)$ is a nontrivial factor of $N$.
\item If $s=\pm r$, retry with a fresh $r$. Each try factors with probability $\tfrac12$, so after $k$ tries the failure probability is $2^{-k}$ -- negligible in a polynomial number of tries.
\end{enumerate}
Run it on $N=3233=61\cdot 53$. Say $r=2$, so $y = 4$. The four roots of $4$ modulo $3233$ are $2,\,1588,\,1645,\,3231$ (that is, $\pm 2$ and $\pm 1645$, arising from the four sign choices in Figure~\ref{ch08:fig-rabin-crt}). Suppose $\mathcal A$ hands back $s = 1645$, which is $\ne \pm 2$: it is congruent to $-2\pmod{61}$ and $+2\pmod{53}$, so $s-r = 1643$ is divisible by $53$ but not by $61$. Then $\gcd(s-2,\,3233)=\gcd(1643,\,3233)=53$, a nontrivial factor of $N$. We have factored $N$ using one call to the root-finder. \emph{Breaking $X$ (taking square roots) solved $Y$ (factoring).}
\end{example}

\begin{figure}[H]
\centering
\begin{tikzpicture}[font=\small, >=Stealth]
  \node[draw=sourceblue!80, fill=sourceblue!10, thick, rounded corners=4pt, minimum width=3.4cm, minimum height=1.2cm, align=center] (sq) at (0, 0) {
    \textbf{Square $y \equiv r^2 \pmod N$}\\[2pt]
    \footnotesize $y = 4 \pmod{3233}$
  };

  % 4 CRT roots
  \node[draw=gray!60, fill=gray!8, thick, rounded corners=3pt, minimum width=3.4cm, minimum height=0.8cm, align=center] (r1) at (5.2, 1.5) {
    \footnotesize $(+r, +r) \mapsto s = 2$ \quad (trivial)
  };
  \node[draw=gray!60, fill=gray!8, thick, rounded corners=3pt, minimum width=3.4cm, minimum height=0.8cm, align=center] (r2) at (5.2, 0.5) {
    \footnotesize $(-r, -r) \mapsto s = 3231$ \quad (trivial)
  };
  \node[draw=targetgreen!80, fill=targetgreen!12, thick, rounded corners=3pt, minimum width=3.4cm, minimum height=0.8cm, align=center] (r3) at (5.2, -0.5) {
    \footnotesize $(-r, +r) \mapsto s = 1645$ \quad \textbf{splits $N$}
  };
  \node[draw=targetgreen!80, fill=targetgreen!12, thick, rounded corners=3pt, minimum width=3.4cm, minimum height=0.8cm, align=center] (r4) at (5.2, -1.5) {
    \footnotesize $(+r, -r) \mapsto s = 1588$ \quad \textbf{splits $N$}
  };

  \node[draw=quotientorange!80, fill=quotientorange!12, thick, rounded corners=4pt, minimum width=3.2cm, minimum height=1.3cm, align=center] (gcd) at (10.2, -1.0) {
    \textbf{$\gcd(s - r, N)$}\\[2pt]
    \footnotesize $\gcd(1643, 3233) = 53$\\[1pt]
    \scriptsize Recovers $p, q$ ($50\%$ chance)
  };

  \draw[->, thick, gray!70] (sq.east) -- (r1.west);
  \draw[->, thick, gray!70] (sq.east) -- (r2.west);
  \draw[->, thick, targetgreen] (sq.east) -- (r3.west);
  \draw[->, thick, targetgreen] (sq.east) -- (r4.west);

  \draw[->, thick, targetgreen] (r3.east) -- (gcd.north west);
  \draw[->, thick, targetgreen] (r4.east) -- (gcd.south west);
\end{tikzpicture}
\caption{Why finding square roots modulo $N = pq$ factors $N$ (Example~\ref{ch08:ex-rabin}). By the Chinese Remainder Theorem, every quadratic residue $y \equiv r^2 \pmod N$ has four roots corresponding to $(\pm r \bmod p,\; \pm r \bmod q)$. Because $y$ hides which root $r$ the reduction chose, any root-finding oracle returns a mixed-sign root $s \not\equiv \pm r \pmod N$ with probability $1/2$, and $\gcd(s-r, N)$ immediately exposes a prime factor.}
\label{ch08:fig-rabin-crt}
\end{figure}

\begin{whybox}
\textbf{What the reduction did and did not prove.} It proved that root-finding mod $N$ is \emph{at least as hard} as factoring: a root-finder yields a factorizer of comparable efficiency. It did \emph{not} prove factoring is hard in any absolute sense -- that remains a conjecture. Reductions move hardness around; they relate problems. The foundation -- ``factoring is hard'' -- is an assumption we adopt because the world's best mathematicians have failed to refute it for half a century. A good security proof is therefore a chain: \emph{scheme} reduces to \emph{root-finding} reduces to \emph{factoring}, and the last link rests on the assumption that factoring is hard.
\end{whybox}

\begin{focusbox}
\textbf{Tightness matters.} A reduction is only as useful as it is efficient. If breaking $X$ with advantage $\eps$ yields a $Y$-solver with success $\eps/q^{2}$ after $q$ queries, then a real attacker with $\eps=1/2$ and $q=2^{40}$ queries gives a $Y$-solver with success only $2^{-81}$ -- weak evidence. A \emph{tight} reduction loses little (constant factors), so breaking $X$ would almost as easily break $Y$. When you read ``provably secure,'' always ask: secure under \emph{which} assumption, and how much does the reduction \emph{lose}? Loose reductions still rule out whole classes of attacks, but they translate into larger recommended key sizes.
\end{focusbox}

\begin{attackbox}
\textbf{When the assumption is the weak link.} A reduction proves ``no efficient attack on $X$ unless $Y$ is easy.'' If $Y$ turns out to be easier than believed, the proof is intact but the scheme is dead. Two real episodes:
\begin{itemize}
\item \emph{Discrete log in weak groups.} Picking the group for a discrete-log scheme carelessly -- a prime $p$ where $p-1$ has only small prime factors -- lets the Pohlig--Hellman algorithm solve discrete log fast. The reduction ``scheme is secure if discrete log is hard'' is fine; the chosen instance simply made discrete log easy. The fix is to choose ``safe'' parameters where the assumption genuinely holds.
\item \emph{Quantum computers.} Shor's algorithm solves both factoring and discrete log in polynomial time on a large quantum computer. Every reduction to those problems survives word-for-word, yet the schemes break, because the assumption ``factoring is hard'' fails against a quantum adversary. This is exactly why post-quantum cryptography (Chapter~35) reaches for \emph{different} hard problems -- lattices, codes -- whose hardness is believed to persist against quantum machines.
\end{itemize}
A proof of security is conditional, so check that its assumption actually holds for the parameters and adversaries you face.
\end{attackbox}

\section{Standard hardness assumptions}
\label{ch08:landscape}

The chapters ahead build schemes on a short list of conjectured-hard problems. It helps to see them gathered, with the one-line statement of each and where it shows up.

\begin{definition}[Foundational hardness assumptions]
\label{ch08:def-assumptions}
\leavevmode
\begin{itemize}
\item \textbf{Factoring.} Given $N=pq$ for unknown primes $p,q$, find $p$ and $q$. Best known classical method: the general number field sieve, sub-exponential time. \emph{Powers:} RSA, Rabin.
\item \textbf{RSA problem.} Given $(N,e,y)$, find $x$ with $x^{e}\equiv y \pmod N$ -- invert the RSA permutation without the factorization. Believed hard; clearly no harder than factoring (factor $N$ and you win), possibly easier (unproven). \emph{Powers:} RSA encryption and signatures.
\item \textbf{Discrete logarithm (DLP).} In a cyclic group $\G=\langle g\rangle$, given $y=g^{x}$, find $x$. \emph{Powers:} Diffie--Hellman, ElGamal, DSA, and -- in elliptic-curve groups -- ECDH and ECDSA.
\item \textbf{Computational Diffie--Hellman (CDH).} Given $g^{a}$ and $g^{b}$, compute $g^{ab}$. No harder than DLP (solve for $a$, then exponentiate). \emph{Powers:} key exchange.
\item \textbf{Decisional Diffie--Hellman (DDH).} Distinguish $(g^a,g^b,g^{ab})$ from $(g^a,g^b,g^c)$ for random $c$. A \emph{decision} problem, strictly easier to attack than CDH in some groups, which is why group choice matters. \emph{Powers:} the security \emph{proofs} of ElGamal and many protocols.
\end{itemize}
\end{definition}

These problems are related by short reductions (Figure~\ref{ch08:fig-assumptions-hierarchy}): $\text{DDH} \le \text{CDH} \le \text{DLP}$, where $\le$ means ``no harder than'' (an algorithm for the right-hand problem yields one for the left). Solving DLP solves CDH solves DDH; the converses are open. Cryptographers like to assume the \emph{weakest} sufficient problem is hard, because a scheme proven secure under a weaker assumption stands on firmer ground.

\begin{figure}[H]
\centering
\begin{tikzpicture}[
  font=\small, >=Stealth,
  prob/.style={draw=#1!75!black, fill=#1!10, thick, rounded corners=4pt, minimum width=3.1cm, minimum height=1.15cm, align=center}
]
  % Factoring lane (top)
  \node[prob=sourceblue] (fact) at (0, 1.3) {
    \textbf{Factoring}\\[1pt]
    \scriptsize Given $N=pq$, find $p,q$
  };
  \node[prob=structurepurple] (rsa) at (5.4, 1.3) {
    \textbf{RSA Problem}\\[1pt]
    \scriptsize Given $(N,e,x^e)$, find $x$
  };
  \draw[->, very thick, sourceblue!80!black] (fact.east) -- node[above=2pt, font=\scriptsize]{compute $d$} (rsa.west);

  % Discrete log lane (bottom)
  \node[prob=sourceblue] (dlp) at (0, -0.5) {
    \textbf{Discrete Log (DLP)}\\[1pt]
    \scriptsize Given $(g,g^a)$, find $a$
  };
  \node[prob=structurepurple] (cdh) at (5.4, -0.5) {
    \textbf{Computational DH}\\[1pt]
    \scriptsize Given $(g^a,g^b)$, find $g^{ab}$
  };
  \node[prob=targetgreen] (ddh) at (10.8, -0.5) {
    \textbf{Decisional DH}\\[1pt]
    \scriptsize Distinguish $g^{ab}$ from $g^c$
  };

  \draw[->, very thick, sourceblue!80!black] (dlp.east) -- node[above=2pt, font=\scriptsize]{$(g^b)^a = g^{ab}$} (cdh.west);
  \draw[->, very thick, structurepurple!80!black] (cdh.east) -- node[above=2pt, font=\scriptsize]{compare $g^{ab}$} (ddh.west);
\end{tikzpicture}
\caption{Reduction hierarchy among standard number-theoretic assumptions. Arrows point from the harder problem to the no-harder problem ($A \to B$ means an oracle for $A$ immediately solves $B$).}
\label{ch08:fig-assumptions-hierarchy}
\end{figure}

\begin{connectionbox}
The two cost claims underpinning these assumptions are developed in the companion books. The algorithms book derives the number field sieve's sub-exponential running time and the $O(\sqrt{|\G|})$ cost of generic discrete-log attacks (baby-step giant-step, Pollard's rho) -- the $\sqrt{}$ being a birthday-bound phenomenon straight out of the probability book. Those two results are why a 3072-bit RSA modulus and a 256-bit elliptic-curve group are reckoned to offer comparable ``128-bit'' security: the elliptic-curve group has no known sub-exponential attack, so it can be far smaller for the same effort.
\end{connectionbox}

\begin{applicationbox}
\textbf{Where this lands in practice.} When your browser opens a TLS connection, the handshake runs an elliptic-curve Diffie--Hellman key exchange -- its secrecy is the CDH/DDH assumption in an EC group, made concrete at curve \texttt{X25519}. The server's certificate is signed with RSA or ECDSA, leaning on factoring or EC discrete log. Your saved passwords sit in a database as one-way hashes (a slow, salted hash like \texttt{argon2} or \texttt{bcrypt}, chosen so that even inverting one digest is costly). WiFi's WPA3 uses a Diffie--Hellman-based handshake (Dragonfly) so that an eavesdropper who records the exchange still faces a discrete-log wall. Signal's double ratchet chains EC Diffie--Hellman steps so that compromising today's key does not unravel yesterday's messages. Every one of these is the same bet -- a problem easy forward, hard backward -- instantiated at a key size chosen so the backward cost dwarfs any attacker's budget. Pick the size from current guidance (a competent library defaults sensibly), not from a textbook's toy numbers.
\end{applicationbox}

\section{Putting the pieces together}
\label{ch08:summary}

The argument of this chapter, end to end: we cannot afford unconditional secrecy, so we settle for security against \emph{efficient} adversaries. ``Efficient'' is polynomial time in the security parameter $n$; the adversary is $\PPT$, randomized and time-bounded. ``Insecure'' means some such adversary wins with \emph{noticeable} probability; ``secure'' means every such adversary's \emph{advantage} -- her edge over guessing -- is \emph{negligible}, shrinking faster than any inverse polynomial as we turn up $n$. The negligible/noticeable line is drawn exactly where it is so that the manipulations of a proof (repeat polynomially often, union-bound, compose) cannot turn safe into unsafe.

On that foundation we built the central object: a \emph{one-way function}, easy forward and hard backward on average, and its richer cousin the \emph{trapdoor} function, hard backward for everyone except the holder of a secret -- the seed of public-key cryptography. We do not know one-way functions exist, but we have well-studied candidates from number theory. Finally, the \emph{reduction} is how we turn a hardness conjecture into a security guarantee: show that breaking the scheme would solve a problem believed hard, and you have proven the scheme secure \emph{relative to} that belief.

Everything after this is specialization. Perfect secrecy (Chapter~9) is the unconditional baseline we are departing from. Stream and block ciphers build fast symmetric mixing. RSA instantiates a trapdoor from factoring; Diffie--Hellman and ElGamal from discrete log; elliptic curves shrink the groups. Signatures, key exchange, TLS, and Signal stack these primitives. When a later proof says ``secure assuming DDH'' or ``a forger would invert RSA,'' you now hold the entire grammar of that sentence: the quantifier over $\PPT$ adversaries, the negligible advantage, the reduction doing the work.

\begin{recoverybox}
Before the exercises, self-check. Can you (1) say why we measure time in the input's \emph{bit-length}; (2) decide whether $n^{-\log n}$ is negligible and justify it; (3) state the two clauses of a one-way function and the loophole each closes; (4) explain what the trapdoor adds and why $pk$ being public is fine; (5) recite the four-step shape of a reduction and what contradiction it reaches? If any of these is fuzzy, reread that section -- the rest of the book assumes all five.
\end{recoverybox}

\section{Exercises}
\label{ch08:exercises}

\begin{exercise}[Negligible or noticeable]
\label{ch08:ex-1}
For each function decide whether it is negligible in $n$ and justify against Definition~\ref{ch08:def-negl}: (a) $2^{-\sqrt n}$; (b) $\dfrac{1}{n^{10}+5n}$; (c) $\dfrac{n^{3}}{2^{n}}$; (d) $\dfrac{1}{(\log n)^{n}}$; (e) $2^{-n} + n^{-2}$. For the noticeable ones, exhibit a constant $c$ for which the function exceeds $1/n^{c}$ infinitely often.
\end{exercise}

\begin{exercise}[Closure properties]
\label{ch08:ex-2}
Prove from the definition that if $\negl_1$ and $\negl_2$ are negligible then so is $\negl_1 + \negl_2$. Then show that for any fixed polynomial $p$, $p(n)\cdot\negl_1(n)$ is negligible. Finally, give a counterexample showing that the \emph{quotient} $\negl_1(n)/\negl_2(n)$ need not be negligible (indeed need not even be bounded).
\end{exercise}

\begin{exercise}[Advantage arithmetic]
\label{ch08:ex-3}
A distinguisher $\mathcal D$ outputs its correct guess of the hidden bit $b$ with probability $\Pr[b'=b] = 0.62$. (a) What is its advantage? (b) A second distinguisher achieves advantage $0.001$; what is its $\Pr[b'=b]$? (c) Suppose a cipher leaks one specific plaintext bit with probability $1/16$ and otherwise looks uniform. Following Example~\ref{ch08:ex-adv}, compute the advantage of the obvious distinguisher.
\end{exercise}

\begin{exercise}[Trace the toy one-way function]
\label{ch08:ex-4}
Work in $\Zns{23}$ with generator $g=5$. (a) Compute $5^{1},\dots,5^{11} \bmod 23$ by hand, each as ``previous $\times 5 \bmod 23$.'' (b) Use your table to solve the discrete log $5^{x}\equiv 9 \pmod{23}$. (c) Roughly how many table entries would the analogous brute-force search scan for a 256-bit modulus, and why does that make the forward/backward asymmetry usable?
\end{exercise}

\begin{exercise}[Break a bad one-way function]
\label{ch08:ex-5}
Consider $f(x) = (3x + 7) \bmod 2^{n}$, claimed by a vendor to be one-way on $n$-bit inputs. Show it is \emph{not} one-way by giving an explicit polynomial-time inverter. (Hint: $3$ is invertible modulo $2^{n}$ because $\gcd(3,2^{n})=1$; find $3^{-1} \bmod 2^{n}$ and solve for $x$.) State the general lesson about why \emph{linear} functions can never be one-way.
\end{exercise}

\begin{exercise}[A reduction by hand]
\label{ch08:ex-6}
Carry out the factoring-from-square-roots reduction of Example~\ref{ch08:ex-rabin} concretely with $N = 77 = 7 \cdot 11$. (a) List all four square roots of $y = 4$ modulo $77$ (use the Chinese Remainder structure: a root is $\pm 2 \bmod 7$ combined with $\pm 2 \bmod 11$). (b) Pick $r = 2$ and a returned root $s \ne \pm 2$; compute $\gcd(s - 2, 77)$ and verify you recover a prime factor. (c) Explain why the reduction fails to learn anything when $\mathcal A$ returns $s = \pm r$, and why retrying fixes this.
\end{exercise}

\begin{exercise}[Order the assumptions]
\label{ch08:ex-7}
Using only the reduction idea ``an algorithm for the harder problem yields one for the easier,'' prove the two relations $\text{CDH} \le \text{DLP}$ and $\text{DDH} \le \text{CDH}$. That is: (a) show a discrete-log oracle lets you solve CDH; (b) show a CDH oracle lets you solve DDH. Conclude that a scheme proven secure under DDH is on weaker (hence safer) ground than one needing DLP to be hard.
\end{exercise}

\begin{exercise}[Why the assumption can fail]
\label{ch08:ex-8}
A protocol designer proves their scheme secure ``assuming discrete log is hard in $\Zns{p}$'' and then picks $p$ such that $p - 1 = 2 \cdot 3 \cdot 5 \cdots$ is a product of small primes. (a) Without computing it in full, explain why Pohlig--Hellman makes discrete log in this group easy, by reducing it to discrete logs in the small prime-order subgroups. (b) Is the security \emph{proof} wrong? Argue precisely what broke -- the proof or the instantiation -- and what ``safe prime'' choice repairs it.
\end{exercise}

\begin{exercise}[Trapdoor, conceptually]
\label{ch08:ex-9}
Explain in your own words, with the padlock analogy, why a \emph{plain} one-way function is enough to store passwords safely but \emph{not} enough to let a stranger send you a private message, whereas a \emph{trapdoor} one-way function enables both. Then identify, for RSA's trapdoor $f(x)=x^e \bmod N$, exactly which piece of data is the public index $pk$ and which is the trapdoor $td$, and why publishing $pk$ does not help an attacker.
\end{exercise}

\begin{exercise}[Baby-step giant-step by hand]
\label{ch08:ex-11}
Reusing the method from the lab box, solve $5^{x}\equiv 6 \pmod{23}$ in $\Zns{23}$. (a) Reuse the baby table $\{5^{0},\dots,5^{4}\}$ and the giant stride $15$; tabulate $6\cdot 15^{i}\bmod 23$ until you hit a stored power, and read off $x = it + j$. (b) Verify by computing $5^{x}\bmod 23$. (c) You did at most $5+5 = 10$ multiplications; the brute-force loop would do up to $22$. State the general count, $O(\sqrt m)$ versus $O(m)$, and explain why this forces a $256$-bit group for $128$-bit security.
\end{exercise}

\begin{exercise}[Find the crossover]
\label{ch08:ex-12}
Consider the negligible function $2^{-n}$ against the polynomial $1/n^{5}$. (a) By trial (or a few lines of code) find the smallest integer $n$ at which $2^{-n} < 1/n^{5}$, and confirm the inequality never flips back. (b) Explain why raising the polynomial to $1/n^{50}$ only moves this crossover to a larger $n$ rather than making $2^{-n}$ non-negligible. (c) Contrast with $1/n^{5}$ versus $1/n^{6}$: why do \emph{two} polynomials never settle the question of negligibility between them?
\end{exercise}

\begin{exercise}[Tightness costs key bits]
\label{ch08:ex-10}
A scheme has a reduction showing that an adversary with advantage $\eps$ making $q$ queries yields a factoring algorithm succeeding with probability $\eps / q^{2}$. Suppose a realistic adversary has $\eps = 1$ and $q = 2^{40}$. (a) With what probability does the derived factorizer succeed, and why is that ``weak'' evidence of security at a modulus where factoring costs about $2^{112}$? (b) Argue qualitatively why a \emph{loose} reduction pushes you toward \emph{larger} key sizes than a tight one would, to keep the guaranteed security level fixed.
\end{exercise}
